\section{Related Work}
\label{sec:related}

\paragraph{VLA Robustness Methods.}
Recent work has sought to improve the robustness of vision-language-action models through training-time interventions.
SRPO~\cite{srpo2025} applies self-referential reinforcement learning to fine-tune VLA policies, achieving strong gains on the LIBERO-Plus benchmark.
RobustVLA~\cite{robustvla2025} combines RL post-training with Jacobian regularization to stabilize action predictions under visual perturbations.
DiG-Flow~\cite{digflow2025} introduces discrepancy-guided flow matching, steering the generative process toward more accurate denoising trajectories during training.
SnapFlow~\cite{luan2026snapflow} distills multi-step flow-matching into single-step inference, reducing latency while preserving action quality.
TT-VLA~\cite{liu2026ttvla} adapts VLA parameters at test time via reinforcement learning, blurring the line between training-time and inference-time approaches.
VLA-GSE~\cite{vla_gse_2026} improves generalization through parameter-efficient fine-tuning with gated sparse experts, and Cosmos-Policy~\cite{cosmos_policy_2026} leverages video foundation models for visuomotor control.
These methods are complementary to ACR-VLA: they improve the base policy itself, whereas ACR-VLA operates as a post-hoc inference-time module that attaches to any frozen flow-matching VLA without modifying its weights.
This orthogonality means ACR-VLA can be composed with any of the above methods to provide additional robustness through uncertainty-conditioned correction.

\paragraph{Uncertainty Quantification for Robot Learning.}
Quantifying predictive uncertainty in learned policies is essential for safe deployment.
Deep Ensembles~\cite{lakshminarayanan2017ensembles} and MC-Dropout~\cite{gal2016dropout} are foundational approaches that estimate epistemic uncertainty through multiple forward passes or stochastic weight perturbations, but both produce scalar or loosely structured uncertainty estimates that do not distinguish which action dimensions are uncertain.
In the generative policy setting, Diffusion Policy~\cite{chi2023diffusion_policy} introduced the diffusion-based action generation paradigm that flow-matching VLAs build upon; its stochastic sampling naturally admits multi-sample uncertainty estimation, though the original work does not exploit this for correction.
DiU~\cite{xing2026diu} shows that for flow-matching models, the divergence of the velocity field recovers per-DOF \emph{aleatoric} uncertainty from a single forward pass, providing an efficient uncertainty signal without additional sampling.
ACR-VLA differs along two axes.
First, it estimates \emph{epistemic} uncertainty via multi-sample covariance---a signal that identifies where the model lacks knowledge rather than where the task is inherently ambiguous---which is more relevant for action correction.
Second, it produces a structured 40-dimensional feature vector (per-DOF standard deviations, calibrated conformal radii, and log-volume) rather than a scalar summary, enabling dimension-selective correction as demonstrated in Section~\ref{sec:motivating}.

\paragraph{Conformal Prediction in Robotics.}
Conformal prediction (CP) provides distribution-free coverage guarantees and has seen growing adoption in robotics.
\citet{angelopoulos2023gentle} provide a comprehensive tutorial on the framework.
Conformalized Quantile Regression~\cite{romano2019cqr} extends CP to heteroscedastic settings, while Adaptive Conformal Inference~\cite{gibbs2021aci} maintains coverage under distribution shift through online threshold updates---a property critical for non-stationary robotic environments.
\citet{barber2023conformal} further relax the exchangeability assumption, enabling conformal methods under general temporal dependence.
Braun et al.~\cite{braun2025cgs} introduce Conformalized Gaussian Scoring for multivariate prediction regions, which ACR-VLA builds upon for its per-DOF feature extraction.
In robotics applications, CP-in-the-Loop~\cite{wang2025cploop} integrates conformal regions into trajectory optimization as constraints, and Conformalized Teleoperation~\cite{zhao2024conformalized} uses CP to provide safety guarantees during shared autonomy.
UPS~\cite{yuan2026ups} employs CP for binary policy gating---deciding whether to execute or abstain---while SAFE~\cite{gu2025safe} uses anomaly detection to trigger safe fallback behaviors.
A common thread in these systems is that CP serves as a \emph{binary decision mechanism} (accept/reject, execute/abstain, safe/unsafe).
ACR-VLA departs from this paradigm by using conformal statistics as \emph{continuous input features} to a learned correction module, enabling graded, dimension-selective refinement rather than all-or-nothing decisions.

\paragraph{Inference-Time Action Correction.}
The most closely related line of work addresses action refinement at inference time without retraining the base policy.
A2C2~\cite{a2c2_2025} learns an asynchronous residual correction that runs in parallel with the base policy's action chunk generation, reducing effective latency.
However, A2C2 applies a generic residual without uncertainty conditioning: it has no mechanism to modulate correction magnitude based on per-dimension confidence, which can lead to unnecessary perturbation of already-accurate action components.
ReconVLA~\cite{chen2026reconvla} applies conformal quantile regression to calibrate per-action uncertainty bounds and uses squared Mahalanobis distance for out-of-distribution detection, providing a principled uncertainty quantification framework that operates without retraining but does not modify the VLA's actions.
ACR-VLA addresses limitations of both approaches.
Compared to A2C2, it conditions correction on observation features extracted from the frozen VLA, enabling the refinement module to produce context-aware corrections.
Compared to ReconVLA, it moves beyond uncertainty quantification to active action improvement, using conformal statistics not just for detection but as part of a two-stage select-then-correct pipeline.
